\documentclass[a4paper,5pt]{amsbook}
%%%%%%%%%%%%%%%%%%%%%%%%%%%%%%%%%%%%%%%%%%%%%%%%%%%%%%%%%%%%%%%%%%%%%

\usepackage{booktabs}
\usepackage{graphicx}
\usepackage{multicol}
\usepackage{textcomp}
\usepackage{systeme}
\usepackage{amssymb}
\usepackage[]{amsmath}
\usepackage{subcaption}
\usepackage[inline]{enumitem}
\usepackage{gensymb}

%%%%%%%%%%%%%%%%%%%%%%%%%%%%%%%%%%%%%%%%%%%%%%%%%%%%%%%%%%%%%%

\newcommand{\sen}{\,\mbox{sen}}
\newcommand{\tg}{\,\mbox{tg}\,}
\newcommand{\cosec}{\,\mbox{cosec}\,}
\newcommand{\cotg}{\,\mbox{cotg}\,}
\newcommand{\tr}{\,\mbox{tr}\,}
\newcommand{\ds}{\displaystyle}
\newcommand{\ra}{\rightarrow}

%%%%%%%%%%%%%%%%%%%%%%%%%%%%%%%%%%%%%%%%%%%%%%%%%%%%%%%%%%%%%%%%%%%%%%%%

\setlength{\textwidth}{16cm} %\setlength{\topmargin}{-1.3cm}
\setlength{\textheight}{25cm}
\setlength{\leftmargin}{1.2cm} \setlength{\rightmargin}{1.2cm}
\setlength{\oddsidemargin}{0cm}\setlength{\evensidemargin}{0cm}

%%%%%%%%%%%%%%%%%%%%%%%%%%%%%%%%%%%%%%%%%%%%%%%%%%%%%%%%%%%%%%%%%%%%%%%%

% \renewcommand{\baselinestretch}{1.6}
% \renewcommand{\thefootnote}{\fnsymbol{footnote}}
% \renewcommand{\theequation}{\thesection.\arabic{equation}}
% \setlength{\voffset}{-50pt}
% \numberwithin{equation}{chapter}

%%%%%%%%%%%%%%%%%%%%%%%%%%%%%%%%%%%%%%%%%%%%%%%%%%%%%%%%%%%%%%%%%%%%%%%

\begin{document}
\thispagestyle{empty}
\pagestyle{empty}
\begin{minipage}[h]{0.14\textwidth}
	\includegraphics[scale=0.24]{../../ufgd.png}
\end{minipage}
\begin{minipage}[h]{\textwidth}
\begin{tabular}{c}
{{\bf UNIVERSIDADE FEDERAL DA GRANDE DOURADOS}}\\
{{\bf C\'alculo Diferencial e Integral III --- Lista 8}}\\
{{\bf Prof.\ Adriano Barbosa}}\\
\end{tabular}
\vspace{-0.45cm}
%
\end{minipage}

%------------------------

\vspace{1cm}
%%%%%%%%%%%%%%%%%%%%%%%%%%%%%%%%   formulario  inicio  %%%%%%%%%%%%%%%%%%%%%%%%%%%%%%%%
\begin{enumerate}
    \setlength\itemsep{0.5cm}
    \item Calcule as integrais iteradas.
        \begin{enumerate}
            \setlength\itemsep{0.3cm}
            \item $\ds\int_0^4\int_0^{\sqrt{y}} xy^2\ dxdy$
            \item $\ds\int_0^1\int_{x^2}^x 1+2y\ dydx$
            \item $\ds\int_0^1\int_0^{x^2} \cos{(x^2)}\ dydx$
        \end{enumerate}

    \item Calcule a integral dupla.
        \begin{enumerate}
            \setlength\itemsep{0.3cm}
            \item $\ds\iint_D \sen{(y^2)}\ dA$, $D=\{(x,y)\ |\ 0\le x\le 1, x\le
            y\le 1\}$
            \item $\ds\iint_D \sen{(y^2)}\ dA$, $D=\{(x,y)\ |\ 0\le y\le 1,
            0\le x\le y\}$
            \item $\ds\iint_D x^2+y^2\ dA$, $D=\{(x,y)\ |\ 0\le x\le 2, x^2\le
            y\le 2x\}$
            \item $\ds\iint_D x^2\ dA$, $D=\{(x,y)\ |\ 1\le x\le e, 0\le y\le
            \ln{x}\}$
        \end{enumerate}

    \item Expresse $D$ como uma regi\~ao do tipo I e tamb\'em do tipo II. Calcule a
    integral dupla das duas formas.
        \begin{enumerate}
            \setlength\itemsep{0.3cm}
            \item $\ds\iint_D x\ dA$, $D$ \'e limitada pelas retas $y=x$, $y=0$ e
            $x=1$.
            \item $\ds\iint_D xy\ dA$, $D$ \'e limitada pelas curvas $y=x^2$,
            $y=3x$.
        \end{enumerate}
\end{enumerate}

\end{document}
